\documentclass[tikz,border=4pt]{standalone}
\usepackage{ctex}
\usepackage{amsmath,amssymb}
\usetikzlibrary{calc,arrows.meta,decorations.markings,patterns}
\begin{document}
\begin{tikzpicture}[
  >=Stealth,
  font=\small,
  ccw/.style={->, very thick, decoration={markings,
      mark=at position 0.3 with {\arrow{>}},
      mark=at position 0.8 with {\arrow{>}}},
    postaction={decorate}},
  cw/.style={->, very thick, decoration={markings,
      mark=at position 0.3 with {\arrow{<}},
      mark=at position 0.8 with {\arrow{<}}},
    postaction={decorate}}
]
\draw[->] (-3,0) -- (4,0) node[right] {$x$};
\draw[->] (0,-3) -- (0,4) node[above] {$y$};

% 外部大区域 C（不规则形）
\fill[blue!8]
  plot[domain=0:360, samples=80, smooth]
    ({2.4*cos(\x)+0.4*cos(2*\x)},
     {2.0*sin(\x)+0.3*sin(2*\x)});

% 挖掉小洞（圆形，含奇点）
\fill[white] (0,0) circle (0.7);
\draw[red!70!black, dashed] (0,0) circle (0.7);

% 奇点
\fill[red] (0,0) circle (2pt);
\node[red, font=\footnotesize, anchor=west] at (0.1, -0.3) {奇点 $(0,0)$};

% 外边界 C（逆时针 = 正向）
\draw[ccw, blue!70!black]
  plot[domain=0:360, samples=80, smooth]
    ({2.4*cos(\x)+0.4*cos(2*\x)},
     {2.0*sin(\x)+0.3*sin(2*\x)});
\node[blue!70!black] at (3.0, 2.0) {$C$ (正向)};

% 内小圆 ℓ 顺时针（使环形区域 D 在左侧）
\draw[cw, orange!80!black] (0,0) circle (0.7);
\node[orange!80!black, anchor=west, font=\footnotesize] at (0.8, 0.6) {$\ell^-$ (顺时针)};

% 环形区域 D 标注
\node[blue!50!black] at (-1.6, -1.5) {$D$ (环形)};

% 公式标注（左上角空白）
\node[font=\small, draw=gray!60, fill=gray!10, fill opacity=0.9,
      text opacity=1, rounded corners=2pt, align=left]
  at (-2.8, 3.8)
  {\textbf{挖洞法}\\[2pt]
   若 $\partial_x Q=\partial_y P$ 在 $D$ 上成立：\\[1pt]
   $\displaystyle\oint_C\!\!=\oint_\ell$（同向）\\[1pt]
   \footnotesize 即使 $C$ 变形，只要绕同样奇点，积分不变};

% 右下说明
\node[font=\footnotesize, align=center] at (1.0, -2.8)
  {环路积分 $\propto$ 被绕奇点数（留数定理雏形）};
\end{tikzpicture}
\end{document}
